\begin{table}[H]
\centering\footnotesize
\caption{Nhiễu iid $N(0,1)$, $n=200$, $B=2000$: tần suất vượt dải và phân phối của $Q_{LB}(10)$}\label{tab:tab_b11_04_simulation}
\begin{adjustbox}{max width=\linewidth}
\begin{tabular}{|>{\raggedright\arraybackslash}p{\dimexpr 0.410\linewidth-2\tabcolsep\relax}|r|r|r|r|}
\toprule
Đại lượng & Mô phỏng & Lý thuyết (tiệm cận) & Xấp xỉ mẫu hữu hạn & MCSE \\
\midrule
Tỷ lệ \textbar{}$\hat{\rho}$(h)\textbar{} vượt $\pm$1.96/$\sqrt{n}$ (h = 1..20) & 0.0435 & 0.05 & 0.0436 & 0.0011 \\[3pt]
Tỷ lệ $Q_{LB}$(10) $>$ phân vị 95\% của $\chi^{2}_{10}$ & 0.0615 & 0.05 & – & 0.0049 \\[3pt]
Trung bình của $Q_{LB}$ & 10.08 & 10 & – & 0.1058 \\[3pt]
Phương sai của $Q_{LB}$ & 22.41 & 20 & – & – \\[3pt]
p-value kiểm định KS so với $\chi^{2}_{10}$ & 0.4655 & – & – & – \\[3pt]
\bottomrule
\end{tabular}
\end{adjustbox}\par\vspace{2pt}{\footnotesize\textit{Ghi chú:} Xấp xỉ mẫu hữu hạn dùng $\E\hat\rho(h)\approx-1/n$ và $\Var\hat\rho(h)\approx(n-h)/\{n(n+2)\}$, lấy trung bình theo $h=1,\dots,20$. MCSE của tỷ lệ vượt dải coi $20B$ chỉ báo là độc lập.}
\end{table}
